% ----------------------------------------------------------------------
% Capítulo 5: Considerações finais
% ----------------------------------------------------------------------
\chapter{Considerações finais}
\label{cap:consideracoes}

\instrucao{De duas a quatro páginas. Retome cada objetivo específico e diga como foi atingido. Aponte limitações, o impacto esperado do produto e sugestões de continuidade. Não traga dados novos.}

O diagnóstico confirmou que a falta de critérios escritos é o principal obstáculo à gestão dos acervos escolares da rede. O guia elaborado responde a esse problema com textos-modelo e formulários simples, e a validação indicou que as equipes conseguem aplicá-lo sem apoio externo.

A pesquisa se limitou a bibliotecas com profissional fixo. Estudos futuros podem testar o guia em unidades atendidas por professores readaptados e acompanhar, ao longo de dois anos, os efeitos da política sobre o uso do acervo.
